\section{Scope, Corpus, and Qualifications}
\label{app:scope}

This study addresses readers approaching Catholic angelology through
Scripture, the Fathers and the systematic teaching of Thomas Aquinas.
Its question is how created spirits receive being, knowledge, love and
beatitude, and how they participate in divine providence. Its exposition
and reference tables are aids to study, not an official catechism or an
ecclesiastical judgment.

The systematic census covers every article of \emph{Summa theologiae} I
questions 50--64 and 106--114. It records the article's question,
determination, authority and grade; it does not reproduce every objection
and reply. The \emph{Celestial Hierarchy} is expounded through its fifteen
chapters. Other works of Aquinas and the patristic, scholastic, Byzantine,
magisterial and liturgical witnesses enter at the identified loci actually
used. Those selections are not an exhaustive corpus of angelology, and the
Dionysian concordance does not claim to count every citation throughout
Aquinas's works.

\textbf{Grades.} ``Defined'' identifies a proposition taught in the named
conciliar profession or definition. ``Church teaching'' identifies other
magisterial teaching. ``Common teaching'' describes agreement supported by
the witnesses named. ``Thomist position'' attributes a systematic
determination to Aquinas; another school's position is named where it is
examined. ``Disputed'' identifies an actual open disagreement in the
witnesses. ``Pious tradition'' and ``Apocryphal'' identify
their respective devotional or extra-canonical status where used. These
grades are editorial assessments of particular propositions; the cited
source supplies their basis. A question containing several propositions
does not receive one undifferentiated doctrinal rank.

\textbf{Language and witnesses.} Biblical quotations follow the registered
Douay--Rheims/Challoner witness. Vulgate book names and numbering are used;
Hebrew psalm numbering is supplied where different. Aquinas's English is
identified with the English Dominican translation and controlled at selected
Latin loci; Dionysius's English follows John Parker's historical translation.
The source register identifies the translation and the degree of checking
for each witness. A searchable transcription is not a critical edition or
a claim of complete manuscript collation. Greek terms are transliterated;
English order names are reconciled in the terminology table.

\textbf{History and interpretation.} The Dionysian corpus's traditional
attribution and its late fifth-/early sixth-century historical setting are
distinguished where they affect a claim. Patristic interpretations are
presented within Catholic theological tradition. Differences among witnesses
are attributed to the authors who express them. The calendar tables identify
their particular books and do not supply a local ordo for an unspecified
diocese. Official web witnesses and devotional discipline were consulted as
of 29 September 2026.

\textbf{Rights and use.} Historical translations and quoted source texts
retain their recorded public-domain or other status; the project's license
does not relicense them. Modern official texts are cited and briefly
paraphrased. No prayer is newly composed or adapted for recitation. The
treatment of fallen angels supplies no method for diagnosing possession
or conducting deliverance. No independent human theological review,
\emph{nihil obstat}, or \emph{imprimatur} is claimed.
